\section*{Chapter \ref{ch:fm:the-economics-of-a-trading-firm} --- The Economics of a Trading Firm}

\begin{solution}{exo:fm:the-economics-of-a-trading-firm:1}
$2\,436.7+617.0+508.8-769.8-647.4=\$2\,145.3$ million. \omterm{def:fm:the-economics-of-a-trading-firm:ntr}{Volume-driven costs} $769.8+647.4=1\,417.2$ against revenue before other income
$2\,145.3+1\,417.2=3\,562.5$: 39.8\%.
\end{solution}

\begin{solution}{exo:fm:the-economics-of-a-trading-firm:2}
$675/1\,398=48.3\%$ in 2025 and $678/1\,696=40.0\%$ in 2022. Pay fell by \$3 million while revenue fell by \$298 million: most pay does not follow
revenue within a year, so the ratio moves against revenue.
\end{solution}

\begin{solution}{exo:fm:the-economics-of-a-trading-firm:3}
$13\,117/79\,748=16.4\%$, as reported.
\end{solution}

\begin{solution}{exo:fm:the-economics-of-a-trading-firm:4}
Net revenue 75. Pay flexing: $0.7\times75-50=2.5$. All pay fixed: $75-80=-5$, a loss.
\end{solution}

\begin{solution}{exo:fm:the-economics-of-a-trading-firm:5}
Net revenue $+133.0\%$, operating profit $+607.5\%$, ratio 4.57. The operating leverage at the 2019 revenue is $1/x^\ast=1/0.229=4.37$ with the
fitted pay slope (4.59 with pay fixed): the arc ratio over a doubling is close to the point elasticity because costs barely moved.
\end{solution}

\begin{solution}{exo:fm:the-economics-of-a-trading-firm:6}
2025: pay \$528.1 million against \$438.0 million on the line. Check whether pay began to follow revenue (a larger variable award in a strong year)
and whether the firm grew: its headcount rose from about 969 to about 1\,027 between the two filings.
\end{solution}

\begin{solution}{exo:fm:the-economics-of-a-trading-firm:7}
Without 2020 the slope is 0.13 with a standard error of 0.13, against $0.29\pm0.094$: one low-revenue, low-pay year carried the estimate. The
slope is a first reading, and the firm's own split of fixed and variable pay is the better source.
\end{solution}

\begin{solution}{exo:fm:the-economics-of-a-trading-firm:8}
The denominators differ: a market maker's equity is what its owners left in the firm, which finances a much larger balance sheet through brokers
and lenders; a bank segment's equity is an allocation set by the group's capital rules. The years differ too: the market maker's ROE was 17.3\% in
2023. Returns on equity compare businesses only when equity is measured the same way and over a cycle.
\end{solution}

\begin{solution}{pb:fm:the-economics-of-a-trading-firm:1}
\begin{enumerate}
\item Revenue less the costs that grow with the business done; financing of positions and margin grows with the positions, so it is one.
\item \$2\,145.3 million, equal to the firm's adjusted net trading income.
\item Market maker \$1\,205.7 million, 56.2\%; asset manager \$428 million, 32.3\%; bank segment \$17\,574 million, 42.8\%.
\item Market maker \$2.09 million and \$0.51 million per head; asset manager \$0.77 million and \$0.39 million.
\item Label and define it, present the nearest statutory (GAAP) figure with it, and reconcile the two quantitatively.
\item $\Pi=(1-b)\mathrm{Rev}-C^{\mathrm{fix}}$.
\item $x^\ast=1-C^{\mathrm{fix}}/((1-b)\mathrm{Rev})=\Pi/((1-b)\mathrm{Rev})$; the leverage is $(1-b)\mathrm{Rev}/\Pi=1/x^\ast$.
\item Market maker $0.048\pm0.045$; asset manager $0.29\pm0.094$; bank segment 0.20 from three years. Only the asset manager's is distinguishable
from zero, and even it moves to 0.13 without 2020.
\item \$836.8 million, \$510.8 million and \$15\,399.9 million.
\item Part of its pay is fixed: when revenue falls, pay falls by less.
\item Flexing: 59.0\%, 45.6\%, 53.3\%. Held fixed: 56.2\%, 32.3\%, 42.8\%.
\item Flexing: 1.69, 2.19, 1.88. Held fixed: 1.78, 3.10, 2.34.
\item 49.2\%, 34.2\% and 43.7\% of 2025's operating profit.
\item All three survive a fall of a third with pay flexing; the asset manager only because its pay flexes (32.3\% break-even with pay fixed).
\item 21.8\%; break-even fall 22.9\% at the 2019 revenue.
\item Market maker \$52 million per index point (standard error 33), correlation 0.58; asset manager $-\$13$ million (23), correlation $-0.27$.
\item 52.7\% in 2025 and 17.3\% in 2023.
\item The market maker's equity is its owners' capital behind a larger financed balance sheet; the bank segment's is an allocation of group equity
under capital rules.
\item The table of the chapter: 59.0\%/1.69 and 56.2\%/1.78; 45.6\%/2.19 and 32.3\%/3.10; 53.3\%/1.88 and 42.8\%/2.34.
\item Make more of pay variable, so that it moves with revenue, and cut fixed costs: both raise the break-even fall. A thin margin cannot be fixed
by hoping for a volatile year.
\end{enumerate}
\end{solution}

\begin{solution}{iq:fm:the-economics-of-a-trading-firm:1}
Pay over net revenue for the period. Part of pay is fixed, so when revenue falls faster than pay the ratio rises, even as pay falls.
\iqlookfor{the definition and the fixed-part argument.}
\end{solution}

\begin{solution}{iq:fm:the-economics-of-a-trading-firm:2}
Volume rises with volatility, and capture per unit (wider spreads) does too; both multiply into revenue.
\iqlookfor{revenue as volume times capture, and both channels.}
\end{solution}

\begin{solution}{iq:fm:the-economics-of-a-trading-firm:3}
The salaried firm: with equal total costs its operating leverage $\mathrm{Rev}/\Pi$ exceeds $(1-b)\mathrm{Rev}/\Pi$, and its break-even fall is
smaller. The bonus firm survives the larger fall.
\iqlookfor{\cref{prop:fm:the-economics-of-a-trading-firm:fall} and the reciprocal relation.}
\end{solution}

\begin{solution}{iq:fm:the-economics-of-a-trading-firm:4}
Revenue 60: $0.8\times60-50=-2$, a loss; profit before the fall was 30.
\iqlookfor{variable costs scale down with revenue, fixed costs do not.}
\end{solution}

\begin{solution}{iq:fm:the-economics-of-a-trading-firm:5}
Equity is measured differently (owners' capital against an attributed share of group equity); the balance sheet is financed differently; the
year and the cycle differ; capital rules differ in kind (chapter 2 against chapter 5).
\iqlookfor{the denominator, not the numerator.}
\end{solution}

\begin{solution}{iq:fm:the-economics-of-a-trading-firm:6}
Four degrees of freedom: about three standard errors from zero, but one year can move it a lot (dropping 2020 halves it here). Headcount growth, an
acquisition or a change in fee mix can raise pay and revenue together. Check the firm's own split of fixed and variable pay.
\iqlookfor{degrees of freedom, leave-one-out fragility, and confounders.}
\end{solution}
